\section{Implementación de la biblioteca}
\label{sec:implementacion}

\subsection{Atributos de calidad}
\label{sec:atributos_calidad}

Los atributos de calidad los elegiremos basado en la ISO 25010, que define un modelo de calidad de software. Este modelo se compone de dos partes: características y subcaracterísticas. Las características son las propiedades generales del software, mientras que las subcaracterísticas son propiedades más específicas. Las mismas se muestran en la \reffig{iso_25010}.

\defaultfigure{implementacion/iso_25010.png}{Modelo de calidad de software ISO 25010.}{iso_25010}

En este caso, buscamos obtener un software que permita a otros investigadores y desarrolladores trabajar de manera eficiente, intuitiva y directa. Que no requieran un entrenamiento extenso para poder utilizarlo. Por lo tanto, los atributos de calidad que consideraremos son de las características Capacidad de Interaccion, y Adecuación Funcional. Particularmente:

\begin{table}[H]
\centering
\begin{tabular}{c c p{6cm}}
\toprule
\textbf{Requisito} & \textbf{Caracteristica} & \textbf{Descripcion} \\ \hline
Adecuacion Funcional & Pertinencia funcional. & Capacidad del producto software para proporcionar un conjunto de funciones que facilitan la consecución de tareas y objetivos de usuario especificados. \\ \hline
Capacidad de interaccion & Auto-descriptividad. & Capacidad de un producto para presentar la información adecuada, haciendo su uso inmediatamente evidente para el usuario sin interacciones excesivas con el producto u otros recursos (documentación, help desks, otros usuarios...). \\ \hline
\bottomrule
% \hline
\end{tabular}
\caption{Atributos de calidad elegidos}
\label{tab:iso25010}
\end{table}

\subsection{Herramientas existentes}

Una vez entendido el problema a resolver y qué atributos de calidad se quieren priorizar, se procedió a analizar el estado del arte en cuanto a herramientas que permitan realizar \textit{quantization aware training}. ¿Se debe construir una biblioteca desde cero? ¿Es posible utilizar otras bibliotecas como base? ¿Existe una solución completa ya disponible?

Tanto en la industria como en los entornos académicos, existen dos bibliotecas de facto \texttt{Tensorflow} y \texttt{PyTorch}.

Ambas tienen un framework de cuantización, \texttt{PyTorch} lo implementa en la misma libreria, a traves del modulo \texttt{\href{https://docs.pytorch.org/docs/stable/quantization.html}{quantization}} y \texttt{Tensorflow} a traves de una libreria secundaria \texttt{\href{https://www.tensorflow.org/model_optimization}{tensorflow-model-optimization}}. En ambos casos las funciones directamente utilizables tienen soporte nativo para 8 bits, pero no suelen soportar múltiples configuraciones de cuantización, sin embargo ambas soportan definir cuantizadores personalizados e inyectarlos en el modelo a cuantizar.

Debido a que ambas bibliotecas parecen tener un soporte similar, y no hay diferencias significativas en la cantidad de usuarios entre una y la otra. Se eligio tensorflow debido a que en el entorno del laboratorio de sistemas embebidos, en el que se desarrolla este trabajo, se utiliza principalmente esta libreria.

\subsection{Metodología de trabajo}

Para desarrollar la libreria, se eligio trabajar con un repositorio de Git, con un flujo de trabajo basado en \textit{Git Flow}. Este flujo de trabajo permite tener una rama principal estable, y ramas de desarrollo que permiten trabajar en nuevas funcionalidades sin afectar la rama principal. Las ramas de desarrollo se fusionan a la rama principal una vez que se ha completado el desarrollo, la revision del codigo y se han realizado las pruebas unitarias necesarias para comprobar la funcionalidad agregada como su compatibilidad con funcionalidades existentes.

\subsubsection{Pruebas unitarias y de integración}

Para la implementación de la biblioteca se utilizó el framework de testing \texttt{unittest} de Python. Este framework permite definir pruebas unitarias y de integración, y ejecutarlas de manera automática. Las pruebas se definen en archivos separados, y se pueden ejecutar de manera individual o en conjunto.

En particular, se usara la metrica de cobertura de código para evaluar la calidad del código. Esta métrica permite medir el porcentaje de código que se ha ejecutado durante las pruebas, y se puede utilizar para identificar áreas del código que no están siendo probadas. Particularmente, se utilizó la herramienta \href{https://coverage.readthedocs.io/}{\texttt{coverage.py}} para medir la cobertura de código, y se generaron informes en formato HTML para visualizar los resultados.

Se obtuvo una cobertura de código del 96\%, lo que indica que la mayoría del código está siendo probado adecuadamente.

\subsection{Analisis de las capacidades y limitaciones de \texttt{TFMOT}}

El mayor punto de dificultad con la libreria \texttt{TFMOT} es la ergonomia para definir esquemas de cuantización complejos. Tiene soporte nativo para cuantizadores uniformes de 8 bits, tanto con signo como sin signo, para otras configuraciones, es necesario definir y articular varias clases e interfaces, como cuantizadores, definiciones de cuantización, y anotaciones de cuantización. Esta falta de ergonomia dificulta la interaccion de los usuarios con la libreria, ya que no es posible definir y aplicar sencillamente diferentes esquemas de cuantización a diferentes capas. Sumado a esto, la documentación oficial es escasa y no provee ejemplos claros de como utilizar estas funcionalidades avanzadas.

El caso sencillo de cuantización de cuantizar un modelo entero con 8 bits se puede realizar como se muestra en el fragmento de codigo \ref{code:tfmot_quantization_simple}.

\begin{lstlisting}[language=Python,label=code:tfmot_quantization_simple,caption=Uso sencillo de TFMOT para cuantización uniforme de 8 bits.]
quant_aware_model = tfmot.quantization.keras.quantize_model(base_model)
\end{lstlisting}

% \texttt{QTensor} fue diseñado como una capa de abstracción y extensión sobre \texttt{TFMOT}. Esto permite a los desarrolladores aprovechar las funcionalidades existentes de \texttt{TFMOT}. La arquitectura modular de \texttt{QTensor} facilita la incorporación de nuevas características y la adaptación a diferentes casos de uso.

En el caso de querer utilizar una configuracion especifica para una capa, se debe realizar lo que se muestra en el fragmento de codigo \ref{code:tfmot_quantization_custom_layer}

\begin{lstlisting}[language=Python,label=code:tfmot_quantization_custom_layer,caption=Uso de TFMOT para cuantización con configuración personalizada en una capa.]
  quantize_config = MyDenseQuantizeConfig()

  model = quantize_annotate_model(
    keras.Sequential([
      layers.Dense(10, activation='relu', input_shape=(100,)),
      quantize_annotate_layer(
          layers.Dense(2, activation='sigmoid'),
          quantize_config=quantize_config)
    ]))

  quantized_model = quantize_apply(model)
\end{lstlisting}
%   Source https://github.com/tensorflow/model-optimization/blob/e2d531bfa2ce10916a1066c67b1438d17cc67dac/tensorflow_model_optimization/python/core/quantization/keras/quantize.py#L83

Como se muestra en el fragmento de codigo \ref{code:tfmot_quantization_custom_layer}, es necesario definir una clase que implemente la interfaz \texttt{QuantizeConfig}, que define como cuantizar los pesos, las activaciones y las salidas de la capa. Luego, es necesario anotar la capa que se desea cuantizar con la configuración personalizada (generalmente en su definición), y finalmente aplicar la cuantización al modelo completo.

Los cuantizadores existentes en la biblioteca tampoco tienen parámetros entrenables, son fijos, lo cual genera que la adaptación sea solo de los parámetros del modelo, y no de los parámetros del cuantizador, lo cual es una característica deseable en nuestra implementación.

\subsection{Implementación de los cuantizadores}

\subsubsection{Clase \texttt{Quantizer}}

La interfaz \texttt{Quantizer} define las operaciones matemáticas que transforman tensores durante la cuantización. Proporciona dos métodos principales:

\begin{enumerate}
    \item \texttt{build}, que construye las variables de estado necesarias para el algoritmo de cuantización.
    \item \texttt{\_\_call\_\_}, que aplica la operación de cuantización a los tensores de entrada.
\end{enumerate}

Sigue el mismo patrón que las capas de Keras, y permite definir variables de estado que se entrenan junto con el modelo. El método \texttt{build} se llama una vez que se conoce la forma de los tensores de entrada, y permite inicializar las variables de estado. El método \texttt{\_\_call\_\_} se llama cada vez que se aplica la cuantización, y debe implementar la lógica de cuantización, esta es la que se coloca en el grafo computacional, permitiendo la cuantización de los tensores y por ende, esta en el camino del gradiente.

\subsubsection{Modificaciones a la clase \texttt{Quantizer}}
Para la implementación de cuantizadores personalizados con parámetros entrenables y gradientes personalizados, se realizaron algunas modificaciones sobre a la clase \texttt{Quantizer} de \texttt{TFMOT}.

Para permitir definir gradientes personalizados, se define un método adicional \texttt{quantize} decorado con \texttt{@tf.custom\_gradient} y dentro se define la función de cuantización y su gradiente. El método \texttt{\_\_call\_\_} llama a este método para aplicar la cuantización con el gradiente personalizado. Este mecanismo es el que se muestra en el fragmento de código \ref{code:quantizer_custom_gradient}.

\begin{lstlisting}[language=Python,label=code:quantizer_custom_gradient,caption=Definición de un cuantizador con gradiente personalizado.]
  class MyQuantizer(Quantizer):
    def __call__(self, inputs, training, weights, **kwargs):
      return self.quantize(inputs, **kwargs)

    @tf.custom_gradient
    def quantize(self, x, **kwargs):
      # Operacion de cuantizacion
      y = ... # Cuantizacion de x

      def grad(dy):
        # Gradiente personalizado
        return dy * ... # Derivada de la operacion de cuantizacion

      return y, grad
\end{lstlisting}

Para permitir definir parámetros entrenables, se modificó el método \texttt{build} para que las variables de estado se definan como tensores entrenables de Keras (\texttt{tf.Variable}). Esto permite que las variables de estado se entrenen junto con el modelo, y se actualicen durante el proceso de entrenamiento a partir de los gradientes calculados.

Los parámetros entrenables de los cuantizadores a definir requieren ciertas características especiales, por ejemplo los niveles y umbrales de la cuantización deben siempre estar ordenados, el parámetro que define el rango $\alpha$ debe ser siempre positivo, etc. Para asegurar que estas características se mantengan durante el entrenamiento, se implementaron restricciones en los parámetros entrenables utilizando \texttt{tf.keras.constraints.Constraint}, se implementaron restricciones personalizadas en el módulo \texttt{constraints}: \texttt{PositiveConstraint} que asegura que un tensor sea siempre positivo, \texttt{LevelConstraint} que asegura que los niveles sean siempre ordenados, esten dentro del rango de cuantización definido por $\alpha$ y que el primer nivel esta fijado al mínimo del rango $m$, y \texttt{ThresholdConstraint} que asegura lo mismo que \texttt{LevelConstraint} pero fijando el ultimo nivel a el máximo del rango $M$. En el fragmento de código \ref{code:quantizer_constraints} se muestra la implementación de la restricción \texttt{PositiveConstraint}.

\begin{lstlisting}[language=Python,label=code:quantizer_constraints,caption=Definición de restricciones para parámetros entrenables.]
  class PositiveConstraint(tf.keras.constraints.Constraint):
    def __call__(self, w):
        return tf.clip_by_value(w, tf.keras.backend.epsilon(), np.inf)
\end{lstlisting}

Se exponen los inicializadores de los parámetros entrenables en el módulo \texttt{initializers}, para permitir que el usuario pueda elegir como inicializar los parámetros entrenables de los cuantizadores. También se exponen los regularizadores en el módulo \texttt{regularizers}, para permitir que el usuario pueda elegir como regularizar los parámetros entrenables de los cuantizadores.

La definición del metodo \texttt{build} se muestra en el fragmento de código \ref{code:quantizer_build}.

\begin{lstlisting}[language=Python,label=code:quantizer_build,caption=Definición del método build de un cuantizador con parámetros entrenables.]
  class MyQuantizer(Quantizer):
    def build(self, tensor_shape, name, layer):
      alpha = layer.add_weight(
          name.join("_alpha"),
          initializer=self.initializer,
          trainable=True,
          dtype=tf.float32,
          regularizer=self.regularizer,
          constraint=PositiveConstraint(),
        )
        return {"alpha": alpha}
\end{lstlisting}

\subsubsection{Cuantizador uniforme}

La implementación de la función de cuantización uniforme definida en la \refeq{uniform_quantizer} se muestra en el fragmento de código \ref{code:uniform_quantizer}.

\begin{lstlisting}[language=Python,label=code:uniform_quantizer,caption=Definición de la función de cuantización uniforme.]
  def uniform_quantizer(x):
    clipped_x = tf.clip_by_value(x, levels[0], levels[-1])
    q = self.delta() * tf.math.floor(clipped_x / self.delta())
\end{lstlisting}

Donde las variables que no estan definidas explicitamente en este fragmento se definen siguiendo la notación de la \ref{sec:uniform}.

\TODO{Considerar agregar el histograma input/output de la funcion de cuantizacion}

Para implementar los gradientes se utilizan las definiciones de la sección \ref{sec:uniform_quantizer_gradient}, y se utilizan en la función \texttt{grad} dentro del método \texttt{quantize}.

El gradiente de la función de cuantización uniforme con respecto a la entrada $x$, $\frac{\partial q_{u}}{\partial x}$, definida en la \refeq{uniform_quantizer_ste_signed} se implementa como se muestra en el fragmento de código \ref{code:uniform_quantizer_gradient_x}.

\begin{lstlisting}[language=Python,label=code:uniform_quantizer_gradient_x,caption=Definición del gradiente de la función de cuantización uniforme con respecto a la entrada $x$.]
  dq_dx = tf.where(
      tf.logical_and(
          tf.greater_equal(x, levels[0]),
          tf.less_equal(x, levels[-1]),
      ),
      upstream,
      tf.zeros_like(x),
  )
\end{lstlisting}

El gradiente de la función de cuantización uniforme con respecto al parámetro $\alpha$, $\frac{\partial q_{u}}{\partial \alpha}$, definida en la \refeq{uniform_quantizer_dq_dalpha} se implementa como se muestra en el fragmento de código \ref{code:uniform_quantizer_gradient_alpha}. Donde \texttt{upstream} es el gradiente acumulado desde las capas siguientes en el grafo computacional.

\begin{lstlisting}[language=Python,label=code:uniform_quantizer_gradient_alpha,caption=Definición del gradiente de la función de cuantización uniforme con respecto al parámetro $\alpha$.]
  dq_dalpha = tf.reduce_sum(q / alpha * upstream)
\end{lstlisting}

Es importante notar la utilización de \texttt{tf.reduce\_sum} en la implementación del gradiente con respecto a $\alpha$. Esto se debe a que $\alpha$ es un escalar, y la función de cuantización se aplica elemento a elemento sobre el tensor de entrada $x$. Por lo tanto, para calcular el gradiente total con respecto a $\alpha$, es necesario sumar los gradientes parciales de cada elemento del tensor de entrada.

\TODO{Mostrar entrenamiento, gif copado y evolucion comun}

\subsubsection{Cuantizador flexible}

La implementación de la función de cuantización flexible definida en la \refeq{flexible_quantizer} se muestra en el fragmento de código \ref{code:flexible_quantizer}.

\begin{lstlisting}[language=Python,label=code:flexible_quantizer,caption=Definición de la función de cuantización flexible.]
  def flexible_quantizer(x):
    # Quantize the levels using uniform quantization
    qlevels = uniform_quantizer(self.levels)

    # Quantize input
    q = tf.zeros_like(x)
    for i in range(self.thresholds.shape[0] - 1):
        q = tf.where(
            tf.logical_and(
                tf.greater_equal(x, self.thresholds[i]),
                tf.less_equal(x, self.thresholds[i + 1]),
            ),
            qlevels[i],
            q,
        )
\end{lstlisting}

El gradiente de la función de cuantización flexible con respecto a la entrada $x$, $\frac{\partial q_{f}}{\partial x}$, es idéntico al del cuantizador uniforme, al ser un STE. Lo mismo ocurre con el gradiente con respecto a $\alpha$, $\frac{\partial q_{f}}{\partial \alpha}$.

El gradiente de la función de cuantización flexible con respecto a los niveles $l$, $\frac{\partial q_{f}}{\partial l}$, descripta en la \refeq{flexible_quantizer_dq_dl} se implementa como se muestra en el fragmento de código \ref{code:flexible_quantizer_gradient_levels}.

\begin{lstlisting}[language=Python,label=code:flexible_quantizer_gradient_levels,caption=Definición del gradiente de la función de cuantización flexible con respecto a los niveles $l$.]
  bin_indices = tf.searchsorted(
      qlevels, tf.reshape(q, (-1,)), side="left"
  )
  q_one_hot = tf.one_hot(bin_indices, depth=qlevels.shape[0])
  dq_dlevel = tf.matmul(
                  tf.transpose(q_one_hot), tf.reshape(upstream, (-1, 1))
              )
              dq_dlevel = tf.reshape(dq_dlevel, (-1,))
\end{lstlisting}

En este caso, la dimensión de la variable de entrada $x$ no corresponde con la dimensión de la variable de los niveles $l$. Por ende, se utiliza la acumulación de los gradientes, para todos los valores de entrada $x_j$ que caen dentro del intervalo definido por los umbrales $t_i$ y $t_{i+1}$. Es decir, todas aquellas entradas que son cuantizadas al nivel $l_i$. Para esto, se utiliza \texttt{tf.searchsorted} para encontrar el índice del nivel correspondiente a cada valor cuantizado, luego se utiliza \texttt{tf.one\_hot} para crear una matriz one-hot con estos índices, y finalmente se utiliza \texttt{tf.matmul} para multiplicar la matriz one-hot transpuesta por el gradiente acumulado desde las capas siguientes en el grafo computacional.

El gradiente de la función de cuantización flexible con respecto a los umbrales $t$, $\frac{\partial q_{f}}{\partial t}$, descripta en la \refeq{flexible_quantizer_dq_dt} se implementa como se muestra en el fragmento de código \ref{code:flexible_quantizer_gradient_thresholds}.

\begin{lstlisting}[language=Python,label=code:flexible_quantizer_gradient_thresholds,caption=Definición del gradiente de la función de cuantización flexible con respecto a los umbrales $t$.]
  dq_dthresholds = tf.zeros_like(thresholds)
  for i in range(1, self.thresholds.shape[0] - 1):
      delta_y = qlevels[i - 1] - qlevels[i]
      delta_x = thresholds[i + 1] - thresholds[i - 1]

      # Only those associated with the 'x' values that
      # Fall within the range of the two borderline levels
      masked_upstream = tf.where(
          tf.logical_and(
              tf.greater_equal(x, self.thresholds[i - 1]),
              tf.less_equal(x, self.thresholds[i + 1]),
          ),
          upstream,
          tf.zeros_like(x),
      )

      update_value = (
          delta_y / delta_x * tf.reduce_sum(masked_upstream)
      )

      dq_dthresholds = tf.tensor_scatter_nd_update(
          dq_dthresholds, indices=[[i]], updates=[update_value]
      )
\end{lstlisting}

En este caso también se utiliza la acumulación de los gradientes, para todos los valores de entrada $x_j$ pues la dimensión de la variable de entrada $x$ no corresponde con la dimensión de la variable de los umbrales $t$. Para esto, se utiliza \texttt{tf.where} para crear una máscara que seleccione solo aquellos valores de $x$ que caen dentro del rango definido por los umbrales $t_{i-1}$ y $t_{i+1}$, luego se utiliza \texttt{tf.reduce\_sum} para sumar los gradientes correspondientes a estos valores, y finalmente se utiliza \texttt{tf.tensor\_scatter\_nd\_update} para actualizar el valor del gradiente correspondiente al umbral $t_i$.

\TODO{Mostrar entrenamiento, gif copado y evolucion comun}

\subsection{Implementación de la descripción de esquemas de cuantización}

\subsubsection{Clase \texttt{QuantizeConfig}}

Para definir como cuantizar una capa, se define la interfaz \texttt{QuantizeConfig}, que permite definir como cuantizar los pesos, las activaciones y las salidas de la capa.

\begin{lstlisting}[language=Python,label=code:quantize_config,caption=Interfaz QuantizeConfig de TFMOT.]
class MyQuantizer(Quantizer):
  ...

class SomeOtherQuantizer(Quantizer):
  ...

class MyQuantizeConfig(QuantizeConfig):
  def get_weights_and_quantizers(self, layer):
    return [(layer.kernel, MyQuantizer()),
            (layer.bias, SomeOtherQuantizer())]

  def set_quantize_weights(self, layer, quantize_weights):
    layer.kernel = quantize_weights[0]
    layer.bias = quantize_weights[1]

  def get_activations_and_quantizers(self, layer):
    return [(layer.activation, SomeOtherQuantizer())]

  def set_quantize_activations(self, layer, quantize_activations):
    layer.activation = quantize_activations[0]

  def get_output_quantizers(self, layer):
    # No cuantizar la salida
    return []

  def get_config(self):
    return {}

\end{lstlisting}

Esto implica definir para cada configuración de cuantización, una clase que implemente la interfaz \texttt{QuantizeConfig}, y defina los cuantizadores a utilizar para los pesos, las activaciones y las salidas de la capa. Lo cual genera una gran cantidad de código repetitivo, y dificulta la definición de nuevas configuraciones de cuantización, de pruebas y de experimentación.

\subsubsection{Clase \texttt{GenerateConfig}}
\label{sec:generate_config}
Para solucionar este problema, se busca obtener una forma de definir configuraciones en formato JSON, y que la librería se encargue de generar las clases necesarias para aplicar la cuantización. Esto se logra definiendo una clase \texttt{GenerateConfig} que recibe un diccionario con la configuración de cuantización, y hereda de \texttt{QuantizeConfig}, implementando los métodos necesarios para definir los cuantizadores a utilizar.

Para adaptar los métodos de la interfaz \texttt{QuantizeConfig} a esta nueva forma de definir las configuraciones, se hizo provecho de la funcionalidad de Python de inspeccionar los atributos de una clase en tiempo de ejecución, utilizando la función nativa \texttt{getattr}.

Por ejemplo, la configuracion de pesos es manejada por los siguientes metodos:

\begin{lstlisting}[language=Python,label=code:generate_config,caption=Metodos para cuantizar los pesos]
    def get_weights_and_quantizers(
        self, layer: Layer
    ) -> list[Tuple[Tensor, Quantizer]]:
        weights_and_quantizers = []
        for weight_attr, quantizer in self.weights.items():
            weights_and_quantizers.append(
                (get_nested_attribute(layer, weight_attr), quantizer)
            )
        return weights_and_quantizers

    def set_quantize_weights(
        self, layer: Layer, quantize_weights: list[Tensor]
    ):
        for attribute, quantized_weight in zip(
            self.weights.keys(), quantize_weights
        ):
            set_nested_attribute(layer, attribute, quantized_weight)

\end{lstlisting}

De esta manera, la definición de nuevas configuraciones de cuantización, y la experimentación con diferentes esquemas de cuantización, se vuelve mucho más sencilla y rápida. El ejemplo del fragmento \ref{code:quantize_config} se puede ahora definir sencillamente como se muestra en el fragmento de código \ref{code:generate_config_usage}.

\begin{lstlisting}[language=Python,label=code:generate_config_usage,caption=Uso de la clase GenerateConfig para definir una configuración de cuantización.]
  config = GenerateConfig(
    weights={
      "kernel": MyQuantizer(),
      "bias": SomeOtherQuantizer(),
    },
    activations={
      "activation": SomeOtherQuantizer(),
    },
  )
\end{lstlisting}

Este aporte es clave, ya que permite con una unica clase, definir cualquier configuración de cuantización a partir de un diccionario, lo cual facilita la definición de nuevas configuraciones, y la experimentación con diferentes esquemas de cuantización.

\subsubsection{Módulo \texttt{qmodel}}
\label{sec:qmodel}

Previamente, la cuantización de un modelo completo con diferentes configuraciones de cuantización para cada capa.

\begin{lstlisting}[language=Python,label=code:custom_cuantization_layout,caption=Uso de TFMOT para cuantización con configuración personalizada.]
    my_config = GenerateConfig(
      weights={
        "kernel": MyQuantizer(),
        "bias": SomeOtherQuantizer(),
      },
      activations={
        "activation": SomeOtherQuantizer(),
      },
    )

    my_other_config = GenerateConfig(
      weights={
        "kernel": SomeOtherQuantizer(),
        "bias": MyQuantizer(),
      },
      activations={
        "activation": MyQuantizer(),
      },
    )

    layer_1 = quantize_annotate_layer(
        Dense(128, activation="relu", name="hidden"),
        my_config,
    )
    layer_2 = quantize_annotate_layer(
      Dense(64, activation="relu", name="hidden_1"),
      my_other_config,),
    layer_3 = Dense(10, activation="softmax", name="output")
    model = quantize_annotate_model(Sequential([layer_1, layer_2, layer_3]))

    with quantize_scope({"MyQuantizeConfig": MyQuantizeConfig}):
        quant_aware_model = quantize_apply(model)
\end{lstlisting}

Para simplificar este proceso se diseño una interfaz de alto nivel, en el modulo \texttt{qmodel} que define las funciones necesarias para cuantizar cualquier modelo definido en \texttt{TensorFlow} a partir de cualquier descripción de esquema de cuantización.

El ejemplo del fragmento \ref{code:custom_cuantization_layout} se puede ahora definir sencillamente como se muestra en el fragmento de código \ref{code:qmodel_usage}.

\begin{lstlisting}[language=Python,label=code:qmodel_usage,caption=Uso del modulo qmodel para cuantizar un modelo a partir de una descripción de esquema de cuantización.]
model = Sequential([
    Dense(128, activation="relu", name="hidden"),
    Dense(64, activation="relu", name="hidden_1"),
    Dense(10, activation="softmax", name="output")
])
qconfig = {
    "hidden": {
        "weights": {
            "kernel": MyQuantizer(),
            "bias": SomeOtherQuantizer(),
        },
        "activations": {
            "activation": SomeOtherQuantizer(),
        },
    },
    "hidden_1": {
        "weights": {
            "kernel": SomeOtherQuantizer(),
            "bias": MyQuantizer(),
        },
        "activations": {
            "activation": MyQuantizer(),
        },
    },
}
apply_quantization(model, qconfig)
\end{lstlisting}

De esta manera no solo se desacopla la definición del esquema de cuantización de la definición del modelo, sino que también se simplifica la interacción del usuario con la librería, permitiendo definir esquemas de cuantización complejos de manera sencilla y rápida. También se esconde la complejidad de la librería \texttt{TFMOT}, permitiendo a los usuarios enfocarse en la definición del esquema de cuantización y el modelo, sin preocuparse por los detalles de implementación de la cuantización.


\subsection{Implementación de computos de complejidad espacial}

Una utilidad relevante para este trabajo es poder calcular la complejidad espacial de un modelo cuantizado, para esto se implementaron las funciones necesarias para calcular la complejidad espacial de modelos sin cuantizar y cuantizados, tanto para el esquema de cuantización uniforme como para el esquema de cuantización flexible.

Para calcular la complejidad espacial de un modelo cuantizado, se itera sobre las capas del modelo, y se calcula la complejidad espacial de cada capa, sumando los resultados. La complejidad espacial de una capa se calcula como la suma del tamaño de los pesos y las activaciones.


\subsubsection{Complejidad para el modelo sin cuantizar}

Para un modelo sin cuantizar se calcula la complejidad espacial como la suma del tamaño de los pesos.

\begin[lstlisting][language=Python,label=code:complexity_unquantized,caption=Cálculo de la complejidad espacial de un modelo sin cuantizar.]
def compute_space_complexity(model: tf.keras.Layer) -> int:
  return sum(
      np.prod(w.shape) * w.dtype.size for w in model.weights
  ) + sum(
      np.prod(a.shape) * a.dtype.size for a in model.activations
  )
\subsubsection{Complejidad para el cuantizador uniforme}

\subsubsection{Complejidad para el cuantizador flexible}

Huffmann

\subsection{Implementación del modulo de flujo de trabajo}

\subsubsection{Clase \texttt{Stage}}

Clase \texttt{StageMetadata}

\subsubsection{Clase \texttt{Pipeline}}
